
\documentclass[UTF8,12pt,a4paper]{ctexart}
\usepackage[margin=2.2cm]{geometry}
\usepackage{amsmath,amssymb,mathtools}
\usepackage{fontspec}
\usepackage{xcolor}
\usepackage{enumitem}
\usepackage{hyperref}
\usepackage{fancyvrb}
\usepackage{fvextra}
\usepackage{microtype}
\usepackage{titlesec}
\usepackage{fancyhdr}
\usepackage{setspace}

\hypersetup{
  colorlinks=true,
  linkcolor=blue!55!black,
  urlcolor=blue!55!black
}

\setstretch{1.18}
\setlist[itemize]{leftmargin=1.8em,itemsep=0.25em,topsep=0.3em}
\setlist[enumerate]{leftmargin=2.1em,itemsep=0.25em,topsep=0.3em}
\titleformat{\section}{\Large\bfseries}{\thesection}{0.8em}{}
\titleformat{\subsection}{\large\bfseries}{\thesubsection}{0.8em}{}
\titleformat{\subsubsection}{\normalsize\bfseries}{\thesubsubsection}{0.8em}{}
\pagestyle{fancy}
\fancyhf{}
\fancyfoot[C]{\thepage}
\renewcommand{\headrulewidth}{0pt}

\DefineVerbatimEnvironment{GuideVerbatim}{Verbatim}{
  breaklines=true,
  breakanywhere=true,
  fontsize=\small,
  frame=single,
  framesep=3mm,
  rulecolor=\color{gray!35},
  bgcolor=gray!7
}

\newcommand{\GuideInlineCode}[1]{\texttt{\detokenize{#1}}}

\begin{document}

\section*{路线 1：QUBO / Ising 建模}

QUBO / Ising 是七条路线的共同语言。你可以把它理解成：把一个选择题改写成能量函数，让“好答案”变成“低能量状态”。

这一步做不好，后面的退火、QAOA、constrained mixer 都会在错误地形上努力。

\subsection*{先看一个选择题}

假设我们要在三个模型中选一个：

\begin{GuideVerbatim}
x_a = 1 表示选择 model_a
x_b = 1 表示选择 model_b
x_c = 1 表示选择 model_c
\end{GuideVerbatim}

收益分别是 9、6、7。原问题是最大化收益：

\[
\operatorname*{maximize} 9 x_{a} + 6 x_{b} + 7 x_{c}
\]

但 QUBO 通常写成最小化能量，所以先改成：

\[
\operatorname*{minimize} -9 x_{a} - 6 x_{b} - 7 x_{c}
\]

如果只有这一步，最优解会把三个变量都设成 1，因为每选一个都能降低能量。可原问题要求“只能选一个”。于是要加约束罚项。

Exactly-one 约束是：

\[
x_{a} + x_{b} + x_{c} = 1
\]

常见 penalty 是：

\[
A \cdot  (x_{a} + x_{b} + x_{c} - 1)^2
\]

当刚好选一个时，括号里是 0，不罚。当一个都不选或选多个时，括号不为 0，就增加能量。

最终 QUBO 是：

\[
\begin{aligned}
\operatorname*{minimize}\\
-9 x_{a} - 6 x_{b} - 7 x_{c}\\
+ A \cdot  (x_{a} + x_{b} + x_{c} - 1)^2
\end{aligned}
\]

这就是 QUBO 的核心动作：目标函数保留偏好，约束变成罚项。

\subsection*{QUBO 标准形式}

QUBO 的标准形式是：

\[
\begin{aligned}
\operatorname*{minimize} E(x) = c + \sum_{i} q_{i} x_{i} + \sum_{i<j} q_{ij} x_{i} x_{j}\\
x_{i} \in \{0, 1\}
\end{aligned}
\]

里面只有三类东西：

\begin{itemize}
\item 常数项 \texttt{c}
\item 单变量项 \texttt{q\_i x\_i}
\item 两变量交互项 \texttt{q\_ij x\_i x\_j}
\end{itemize}

为什么只允许二次？因为很多退火器、Ising solver、QAOA cost Hamiltonian 都天然吃二次形式。超过二次的项通常要 quadratization，引入辅助变量降阶。

\subsection*{二次项的直觉}

线性项表示单独选择某个变量的好坏。

\[
-9 x_{a}
\]

代表 \texttt{x\_a = 1} 会让能量下降 9，因此倾向选择它。

二次项表示两个变量同时为 1 时的额外效果。

\[
+20 x_{a} x_{b}
\]

代表同时选择 \texttt{a} 和 \texttt{b} 会被重罚。

\[
-3 x_{c} x_{boost}
\]

代表 \texttt{model\_c} 和 \texttt{boost} 同时出现时有协同收益。

所以 QUBO 不只是“选或不选”，它能表达组合效应。

\subsection*{约束为什么平方}

等式约束通常长这样：

\[
g(x) = 0
\]

把它变成 penalty：

\begin{GuideVerbatim}
A * g(x)^2
\end{GuideVerbatim}

平方有两个好处：

\begin{itemize}
\item 不管违反方向是正还是负，都会罚。
\item 满足约束时 penalty 正好为 0。
\end{itemize}

对于 exactly-one：

\[
g(x) = x_{a} + x_{b} + x_{c} - 1
\]

如果选一个，\texttt{g(x) = 0}。如果选两个，\texttt{g(x) = 1}。如果三个都选，\texttt{g(x) = 2}，罚得更重。

\subsection*{Ising 形式}

Ising 变量不是 0/1，而是 spin：

\[
s_{i} \in \{-1, +1\}
\]

Ising 能量写作：

\[
H(s) = C + \sum_{i} h_{i} s_{i} + \sum_{i<j} J_{ij} s_{i} s_{j}
\]

QUBO 和 Ising 可以互相转换。常用关系：

\[
\begin{aligned}
s_{i} = 2 x_{i} - 1\\
x_{i} = (s_{i} + 1) / 2
\end{aligned}
\]

你可以把 QUBO 看成计算机科学味道的写法，把 Ising 看成物理味道的写法。它们表达的是同一个能量地形。

\subsection*{变量编码}

很多原问题变量不是 binary。

例如一个整数变量：

\[
y \in \{0, 1, 2, 3, 4, 5, 6, 7\}
\]

可以用 binary expansion：

\[
y = b_{0} + 2 b_{1} + 4 b_{2}
\]

也可以用 one-hot：

\[
\begin{aligned}
y = 0\cdot z_{0} + 1\cdot z_{1} + ... + 7\cdot z_{7}\\
z_{0} + z_{1} + ... + z_{7} = 1
\end{aligned}
\]

binary expansion 省变量，但约束和目标可能变复杂。one-hot 变量多，但表达互斥和选择结构很直观。比赛中，编码方式经常决定你能不能把模型压到可求规模。

\subsection*{高阶项二次化}

如果目标里出现三变量乘积：

\[
x_{i} x_{j} x_{k}
\]

它不是 QUBO。常见处理是引入辅助变量：

\[
y = x_{i} x_{j}
\]

然后把高阶项改成：

\[
y x_{k}
\]

再加 penalty 保证 \texttt{y} 真的等于 \texttt{x\_i x\_j}。

直觉上，辅助变量是在帮你记住一个局部组合状态。但它会增加变量数和约束复杂度，不能滥用。

\subsection*{建模流程}

一条健康的 QUBO 建模流程是：

\begin{GuideVerbatim}
1. 定义原问题变量
2. 选择 binary / one-hot / domain-wall / integer encoding
3. 写原始目标函数
4. 把 maximization 改成 minimization
5. 把硬约束变成 penalty 或交给其他机制
6. 展开成线性项和二次项
7. 检查系数范围、变量数、coupler 数
8. 小规模 brute force 验证
9. 解码回原问题变量
10. 重新计算原始目标和约束
\end{GuideVerbatim}

第 10 步非常重要。不要把 QUBO energy 当成业务目标。QUBO energy 包含 penalty 和 offset，它只是求解器看到的地形。

\subsection*{比赛时能改哪些地方}

变量设计

同一个问题可以有多种变量。路径问题可以用边变量，也可以用时间-位置变量。排产问题可以用开始时间变量，也可以用时间槽 one-hot。变量一变，整个 QUBO 难度都变。

编码方式

binary、one-hot、domain-wall 的变量数、耦合密度、可行性控制都不同。比赛时很容易靠编码改进拿到效果。

Penalty 权重

\texttt{A} 太小，非法解会赢。\texttt{A} 太大，目标函数差异被淹没，求解器只看到约束墙。最好做 penalty sweep。

缩放

退火器和量子线路都对系数范围敏感。系数跨度太大，会让小差异不可见。

拆分

大问题不一定直接全量 QUBO。可以固定一部分变量，把剩下的不确定变量变成小 QUBO。

\subsection*{常见坑}

把最大化忘记取负。

这会让算法找到最差解。

Penalty 没有压过目标收益。

例如违反约束能多拿 100 分，但 penalty 只有 20，求解器当然会违法。

只看 raw energy。

最低 QUBO energy 可能不可行。最后答案必须经过 decode、feasibility check、objective recomputation。

变量爆炸。

一个看似自然的 one-hot 编码可能让变量数从几十变成几千。要随时估算规模。

高阶项降阶过度。

每个辅助变量都要约束它的语义。辅助变量太多，模型会变硬。

\subsection*{和其他路线的关系}

路线 2 接着处理约束质量。

路线 3 直接拿 QUBO/BQM 退火采样。

路线 4 把 QUBO/Ising 变成 cost Hamiltonian。

路线 5 会问：有些约束是否不要 penalty，而让 mixer 保持？

路线 6 会问：哪些变量不用全进 QUBO，可以经典求解或固定？

路线 7 会问：QUBO 图里哪些变量更重要，哪些变量能被学习策略 warm-start 或高置信度固定？

\subsection*{共读任务}

用 \texttt{sample\_problem.json} 手算一个简化 QUBO：

\begin{enumerate}
\item 只保留 \texttt{model\_a}、\texttt{model\_b}、\texttt{model\_c}。
\item 写出收益项。
\item 加上 exactly-one penalty。
\item 比较 \texttt{100}、\texttt{010}、\texttt{001}、\texttt{110} 的能量。
\end{enumerate}

如果你能解释为什么 \texttt{110} 会被罚，路线 1 就入门了。

\subsection*{延伸阅读}

\begin{itemize}
\item \texttt{literature/requirements/01\_qubo\_ising\_modeling\_requirements.md}
\item Glover, Kochenberger, Du, QUBO tutorial
\item Lucas, Ising formulations of many NP problems
\end{itemize}

\end{document}
